% !TEX root = ../main.tex
%==============================================================
\chapter{Giới thiệu đề tài}
\label{ch:gioi-thieu}

\section{Đặt vấn đề}
Kết xuất đồ họa (rendering) là quá trình tạo ra hình ảnh hai chiều từ mô tả
của một cảnh ba chiều. Trong số các kỹ thuật kết xuất, \textit{Path Tracing}
là phương pháp dựa trên mô phỏng vật lý của ánh sáng, cho phép tạo ra hình ảnh
có độ chân thực cao với các hiệu ứng như phản xạ, khúc xạ, đổ bóng mềm và
chiếu sáng toàn cục (global illumination). Tuy nhiên, chi phí tính toán của
Path Tracing rất lớn do phải mô phỏng hàng triệu tia sáng cho mỗi khung hình.

% TODO: Bổ sung dẫn nhập về nhu cầu kết xuất thời gian thực.

\section{Mục tiêu của đề tài}
Đề tài hướng tới việc nghiên cứu và xây dựng một bộ kết xuất Path Tracing
chạy song song trên kiến trúc CUDA của NVIDIA nhằm đạt tốc độ tương tác
(thời gian thực). Các mục tiêu cụ thể bao gồm:
\begin{itemize}[noitemsep]
  \item Hiện thực thuật toán Path Tracing trên GPU bằng CUDA.
  \item Xây dựng hệ thống vật liệu Lambertian, Metal, Dielectric và camera vật lý
        có độ sâu trường ảnh, nhòe chuyển động.
  \item Hiển thị thời gian thực qua cầu nối CUDA--OpenGL và tích lũy mẫu theo
        thời gian để khử nhiễu.
  \item Nạp cảnh theo định dạng Mitsuba 3 và cung cấp giao diện điều khiển tương tác.
\end{itemize}

\section{Phạm vi và đối tượng nghiên cứu}
% TODO: Mô tả phạm vi (loại cảnh hỗ trợ, định dạng nhập, phần cứng thử nghiệm).

\section{Cấu trúc báo cáo}
Báo cáo được tổ chức thành năm chương. Chương~\ref{ch:gioi-thieu} giới thiệu đề tài.
Chương~\ref{ch:ly-thuyet} trình bày nền tảng lý thuyết. Chương~\ref{ch:trien-khai}
mô tả quá trình triển khai hệ thống. Chương~\ref{ch:danh-gia} trình bày kết quả
và đánh giá. Chương~\ref{ch:ket-luan} đưa ra kết luận và hướng phát triển.
